% Options for packages loaded elsewhere
\PassOptionsToPackage{unicode}{hyperref}
\PassOptionsToPackage{hyphens}{url}
\PassOptionsToPackage{dvipsnames,svgnames,x11names}{xcolor}
\documentclass[
  11pt,
]{article}
\usepackage{xcolor}
\usepackage[margin=2.4cm]{geometry}
\usepackage{amsmath,amssymb}
\setcounter{secnumdepth}{-\maxdimen} % remove section numbering
\usepackage{iftex}
\ifPDFTeX
  \usepackage[T1]{fontenc}
  \usepackage[utf8]{inputenc}
  \usepackage{textcomp} % provide euro and other symbols
\else % if luatex or xetex
  \usepackage{unicode-math} % this also loads fontspec
  \defaultfontfeatures{Scale=MatchLowercase}
  \defaultfontfeatures[\rmfamily]{Ligatures=TeX,Scale=1}
\fi
\usepackage{lmodern}
\ifPDFTeX\else
  % xetex/luatex font selection
\fi
% Use upquote if available, for straight quotes in verbatim environments
\IfFileExists{upquote.sty}{\usepackage{upquote}}{}
\IfFileExists{microtype.sty}{% use microtype if available
  \usepackage[]{microtype}
  \UseMicrotypeSet[protrusion]{basicmath} % disable protrusion for tt fonts
}{}
\usepackage{setspace}
\makeatletter
\@ifundefined{KOMAClassName}{% if non-KOMA class
  \IfFileExists{parskip.sty}{%
    \usepackage{parskip}
  }{% else
    \setlength{\parindent}{0pt}
    \setlength{\parskip}{6pt plus 2pt minus 1pt}}
}{% if KOMA class
  \KOMAoptions{parskip=half}}
\makeatother
\setlength{\emergencystretch}{3em} % prevent overfull lines
\providecommand{\tightlist}{%
  \setlength{\itemsep}{0pt}\setlength{\parskip}{0pt}}
% v3 header: single column, tight reference list, clickable links.
% Loaded via pandoc -H. Fixes for the six issues Havid found in the v2 PDF are
% mostly in the markdown generator; this file only carries LaTeX-side settings.
\usepackage{graphicx}
\usepackage{float}
\usepackage{caption}

% Figure captions: LaTeX owns the "Figure N." label. The v2 PDF rendered
% "Figure 4: Figure 4." because the caption text ALSO began with "Figure 4.",
% written into the markdown by the generator. v3 passes a bare caption and lets
% LaTeX number it, so the label appears exactly once.
\captionsetup[figure]{labelfont=bf, labelsep=period, font=small,
                      justification=justified, singlelinecheck=false,
                      skip=6pt}

% Tighten the numbered reference list produced by pandoc's default
% CSLReferences environment.
\usepackage{etoolbox}
\AtBeginEnvironment{CSLReferences}{%
  \small
  \setlength{\parskip}{0pt}%
  \setlength{\itemsep}{0pt}%
  \setlength{\topsep}{2pt}%
}

% No section numbering from LaTeX: the generator already writes "1. Introduction",
% "2. The Certified Frontier ...", so LaTeX numbering would double it.
\setcounter{secnumdepth}{0}
\usepackage{bookmark}
\IfFileExists{xurl.sty}{\usepackage{xurl}}{} % add URL line breaks if available
\urlstyle{same}
% fallback for those not using the hyperref driver hyperxmp:
\makeatletter
\@ifundefined{xmpquote}{\newcommand{\xmpquote}[1]{#1}}{}
\makeatother
\hypersetup{
  colorlinks=true,
  linkcolor={[HTML]{1A4E8A}},
  filecolor={Maroon},
  citecolor={Blue},
  urlcolor={[HTML]{1A4E8A}},
  pdfcreator={LaTeX via pandoc}}

\author{}
\date{}

\begin{document}

\setstretch{1.05}
\begin{center}
{\LARGE\bfseries Perovskite Photovoltaics in August 2026: Buried-Interface Chemistry, Defect Tolerance and the Certified Efficiency Frontier\par}
\vspace{1.1em}
{\large Havid Aqoma\textsuperscript{1,2,3,4}\par}
\vspace{0.7em}
\begin{minipage}{\textwidth}\centering\small
\textsuperscript{1}School of Energy and Chemical Engineering, Xiamen University Malaysia, Selangor Darul Ehsan 43900, Malaysia\\[0.25em]
\textsuperscript{2}Kelip-kelip! Center of Excellence for Light Enabling Technologies, School of Energy and Chemical Engineering, Xiamen University Malaysia, Selangor Darul Ehsan 43900, Malaysia\\[0.25em]
\textsuperscript{3}Centre of Excellence for Green and Advanced Technologies in Efficient Separations (GATES), School of Energy and Chemical Engineering, Xiamen University Malaysia, Selangor Darul Ehsan 43900, Malaysia\\[0.25em]
\textsuperscript{4}Gulei Innovation Institute, Xiamen University, Zhangzhou 363200, China\\[0.25em]
\end{minipage}
\vspace{2.2em}
{\small September 8, 2026\par}
\end{center}
\vspace{1.0em}

\subsection{Abstract}\label{abstract}

Perovskite photovoltaics published 531 indexed works in August 2026, of
which 144 were read closely and 82 are cited here. This review is
organised mechanistically rather than chronologically, and separates
what the month's literature verified from what it claimed. The certified
frontier advanced on three fronts, each through interfacial rather than
bulk chemistry. A non-anchoring additive that suppresses self-assembled
monolayer aggregation while leaving substrate binding sites available
delivered a 33.66\% certified perovskite/silicon tandem {[}1{]}; a
thienyl self-assembled monolayer templating preferential crystallisation
of the wide-bandgap subcell reached 32.56\% certified on a tunnel-oxide
passivated contact silicon bottom cell {[}2{]}; and a thiocyanate
healing treatment driving secondary crystal growth into conformal
columnar grains gave 30.77\% certified on industrial textured silicon
{[}3{]}. Single-junction certification clustered near 27.4\%, led by
redirection of a wet-interfacial iodide redox pathway at an acidic
phosphonic-acid contact at 27.4\% certified {[}4{]}. The largest
certified area was a 61.2 cm\(^2\) module at 17.35\% certified, coated
by a closed-loop slit-flushing antisolvent route in ambient air {[}5{]},
which is also the month's clearest statement of the area penalty. Only
18 of 144 closely read papers reported an independently certified power
conversion efficiency at all. Operational-stability evidence was the
weakest link: 5 papers reported a time to 80\% of initial performance
(T80), the longest at 1400 h {[}6{]}, and 2 named a protocol from the
International Summit on Organic Photovoltaic Stability (ISOS), so most
lifetime claims this month cannot be compared against each other. Across
the corpus an efficiency value was detected in at least 55.5\% of
abstracts, independent certification in at least 9.7\%, and a complete
open-circuit voltage, short-circuit current density and fill factor
triplet in at least 4.1\%. Three mechanistic conclusions follow.
Buried-interface chemistry, not absorber composition, now sets the
achievable open-circuit voltage. Defect tolerance has overtaken defect
elimination as the practical design strategy. Halide homogenisation in
wide-bandgap absorbers is the rate-limiting step for tandem integration,
and it is the subcell least often optimised in the month's work.

\textbf{Keywords:} perovskite photovoltaics; buried interface; defect
passivation; wide-bandgap perovskite; tandem solar cells; operational
stability

\subsection{1. Introduction}\label{introduction}

Perovskite single junctions are now certified above 27\% power
conversion efficiency (PCE) and perovskite/silicon tandems above 33\%,
so the field's open problems have moved from raw efficiency to whether a
claimed efficiency is verifiable, whether it survives at module area,
and whether it survives operation. In August 2026, 531 indexed works met
the scope gate and 144 were read closely.

This review follows four mechanistic threads: buried-interface and
contact chemistry, defect tolerance and ion migration, wide-bandgap
absorbers for tandems, and operational stability, with scale-up as the
constraint coupling them. Buried interfaces dictate band alignment and
non-radiative recombination kinetics, while mobile ionic defects screen
internal fields and accelerate chemical corrosion at selective contacts.
In wide-bandgap absorbers required for tandems, light-induced halide
segregation pins photovoltages, whereas rapid crystallization during
large-area coating multiplies carrier-trapping boundaries across the
absorber volume.

The month's headline numbers are strong, the evidence behind them is
unevenly reported, and this review separates the two.

\subsection{2. The Certified Frontier in August
2026}\label{the-certified-frontier-in-august-2026}

Independent verification remains rare across current publications: only
18 of 144 closely read papers reported an independently certified power
conversion efficiency (PCE). For two-terminal (2T) tandems, independent
certification established a 33.66\% ceiling for perovskite/silicon
tandems alongside certified efficiencies of 32.56\% and 30.77\%
{[}1,2,3{]}, whereas all-perovskite tandem devices reached certified
boundaries of 29.7\% and 27.84\% {[}7,8{]}. In single-junction
positive-intrinsic-negative (p-i-n) devices, non-radiative recombination
suppression defines a certified ceiling of 27.4\%, closely corroborated
by certified metrics of 27.2\% and 26.77\% {[}4,9,10{]}.

A self-reported champion on a sub-0.1 cm2 cell and a certified value on
a module are different classes of evidence. Self-reported metrics run
substantially ahead of certified baselines because uncertified testing
on sub-0.1 cm2 apertures artificially minimizes series resistance,
avoids lateral coating non-uniformities, and omits interconnection dead
space. Uncertified claims at larger areas illustrate this tension:
self-reported efficiencies reached 27.7\% on 25 cm2 and 27.01\% on 58.51
cm2 {[}11,12{]}, rivaling certified small-area cells. However,
uncertified performance across similar apertures diverges steeply due to
sheet-resistance and defect-density penalties, falling to 23.3\% on 57.3
cm2, 23.14\% on 29.7 cm2, and 17.35\% on 61.2 cm2 {[}5,13,14{]}. Without
independent verification across macroscopic apertures, uncertified
records mask interfacial heterogeneity and distort true progress.

\begin{figure}[htbp]
\centering
\includegraphics[width=0.95\textwidth]{fig/F1_certified_frontier.pdf}
\caption{The certified frontier in August 2026. Filled markers are independently certified efficiencies, open markers are self-reported champion values, grouped by device family. Horizontal bars mark the highest certified value in each family.}
\end{figure}

\subsection{3. Buried Interfaces and Contact
Chemistry}\label{buried-interfaces-and-contact-chemistry}

Buried contacts in inverted (p-i-n) and regular (n-i-p) perovskite solar
cells (PSCs) govern non-radiative recombination and charge extraction,
yet unpassivated defects, band misalignment, and interfacial voiding
compromise performance {[}15,16,17{]}. At metal-oxide electron-transport
layers (ETLs), reactive surface species trigger mechanical voiding;
gas-solid ethanol vapor reconstruction converts tin oxide (SnO2) surface
hydroxyls into anchored acetate and passivates oxygen vacancies for an
uncertified 26.19\% champion power conversion efficiency (PCE) {[}18{]}.
Subliming volatile pyruvic acid at 145 °C similarly cleans SnO2 to lower
oxygen defects, delivering a 26.17\% champion PCE on a 1 cm2 aperture
with an 80\% initial performance retention (T80) of 900 hours {[}19{]}.
Bidentate bridges such as 3,4-diaminobenzoic acid (DABA) bridge this
boundary by anchoring carboxyl groups to Sn4+ while passivating
undercoordinated lead (Pb2+) defects with amino groups {[}20{]}.

At hole-transport layers (HTLs), interfacial acidity and defect states
provoke analogous degradation. In methylammonium-free lead-tin devices,
4-(trifluoromethyl)benzamidine hydrochloride (TFBA) suppresses
acid-induced degradation by inhibiting sulfonate (-SO3H) deprotonation
on poly(styrenesulfonate) (PSS), securing an uncertified 22.48\%
champion PCE {[}21{]}. On nickel oxide (NiOx), ammonium thiocyanate
(NH4SCN) consumes surface hydroxyls and relieves lattice strain for a
21.41\% champion PCE {[}22{]}. Decoupling self-assembled monolayer (SAM)
chemisorption from bulk transport via a structural-compensation stack of
sputtered NiOx and solution-processed hydroxyl-rich overlayers yields a
24.86\% champion PCE across a 12.4 cm2 aperture {[}23{]}, while
enamine-based SAM anchor engineering achieves 23.14\% PCE and 83.11\%
fill factor (FF) on a 29.7 cm2 mini-module {[}14{]}.

However, SAM hole contacts frequently suffer from aggregation and dipole
misalignment. Non-anchoring 1,3,5-tris(trifluoromethyl)benzene (TTMB)
suppresses SAM aggregation while preserving substrate binding sites,
reducing recombination and shunting to yield certified and stabilised
33.66\% PCEs in perovskite/silicon tandems {[}1{]}. In wide-bandgap
cells, building an interfacial hydrogen-bond network with glycinamide
hydrochloride suppresses molecular clustering in
{[}4-(3,6-dimethoxy-9H-carbazol-9-yl)butyl{]}phosphonic acid (MeO-4PACz)
for a 22.38\% champion PCE {[}24{]}. Asymmetric sulfur- and
chlorine-functionalized molecules tune buried band offsets to achieve
champion PCEs of 26.49\% (25.70\% steady-state) {[}25{]} and 26.38\%
{[}26{]}, while co-assembling dipoles on NiOx yields a 20.36\%
four-terminal tandem champion PCE {[}27{]}. Yet dipolar engineering
often lacks rigorous validation: bismuth oxychloride (BiOCl) flips
interfacial dipoles alongside chlorine-iodine exchange to yield a
certified 27.84\% tandem PCE (champion 28.14\%, stabilised 27.86\%), but
single-junction inverted devices reaching 27.25\% (1.55 eV) and 20.24\%
(1.85 eV) PCE were restricted to an unscaled 0.09 cm2 area {[}8{]},
while dipole tuning with 2-fluorobenzene-1,3-dicarboxylic acid (m-FBC)
reached a 25.03\% champion PCE whose 94\% retention after 500 hours
derived from unlabelled storage {[}28{]}.

Buried contacts also suffer thermomechanical failure under operational
stress. Anchoring {[}60{]}fullerene-phosphonic acid (C60-PA) on SnO2
elevates maximum power point (MPP) tracked efficiency from an average of
18.6\% to 20.0\%, yet degrades under illuminated heating because
fullerenes weaken interfacial adhesion {[}29{]}. Strengthening adhesion
requires molecular templating: pairing precursor additives with
phosphonic acid self-assembled monolayers (SAMs) enhances interfacial
adhesion and suppresses groove depths for a 26.58\% champion PCE
{[}30{]}, while self-assembled multilayers with
{[}2-(10-(3,5-dimethoxyphenyl)-7H-benzo{[}c{]}carbazol-7-yl)ethyl{]}phosphonic
acid (DMPA) secure certified 27.2\% and champion 27.59\% PCEs, retaining
94.5\% efficiency after 1600 hours at 65 °C under 1 Sun illumination
{[}9{]}. Coordinate bonding offers comparable reinforcement:
thienyl-based additives template (111) crystallization and segregate at
grain boundaries to form passivating lead-sulfur (Pb-S) bonds, yielding
a certified 32.56\% (champion 33.0\%) tandem PCE {[}2{]}. Pb-S
coordination via dibutyl sulfide induces a surface p-to-n transition for
an uncertified 26.28\% champion PCE {[}31{]}, while multi-arm
donor-acceptor dipole bridges coordinate Pb2+ via C=N and methoxy groups
to relieve residual strain, yielding a 26.79\% champion PCE {[}32{]}.

\subsection{4. Defect Tolerance, Passivation and Ion
Migration}\label{defect-tolerance-passivation-and-ion-migration}

Undercoordinated lead (Pb2+) sites and halide vacancies drive
non-radiative recombination and initiate photochemical degradation.
Multidentate Lewis bases, chelating additives, and biopolymers suppress
these trap states by coordinating Pb2+ and arresting formamidinium (FA+)
deprotonation or iodide oxidation {[}33,34,35,36,37,38{]}. Coordinating
Pb2+ with phenylthiourea (TPU) retards ion migration by suppressing
iodine vacancies and volatile iodine loss, achieving an uncertified
champion power conversion efficiency (PCE) of 26.13\% {[}33{]}. Ammonium
alginate stabilizes FA+ against deprotonation and passivates Pb2+,
increasing the fill factor (FF) from 80.2\% to 85.1\% in an uncertified
26.12\% cell {[}34{]}. Related treatments using sulfamethoxazole (SMX)
(champion 21.05\%, average 20.31\%) {[}35{]}, a sulfonium-amino acid
additive (champion 22.5\%) {[}36{]}, multi-pyridyl spirobifluorenes
(champion 20.98\% versus 18.67\% for control) {[}37{]}, and
3-(trifluoromethyl)thiophenol (22.30\% PCE) {[}38{]} similarly
neutralize coordination defects, yet these studies examine uncertified
small-area cells without establishing thermal bond durability.

Beyond surface coordination, mitigating ion migration requires
addressing lattice mechanics, strain, and kinetic barriers within the
perovskite absorber {[}39,40,41,42,43,44{]}. Stiffening the framework
through solid-solution-induced lattice hardening via partial lanthanum
substitution for Pb increases the Young's modulus from 27.96 to 35.71
GPa, delivering a certified steady-state efficiency of 26.03\% (champion
26.54\%) {[}39{]}. Regulating crystallographic strain further restricts
ion migration channels: shifting lattice strain from tensile to
compressive with bimolecular sulfonates delivers an uncertified champion
26.21\% PCE {[}40{]}, dynamic strain homogenization via supramolecular
networks in tin-lead (Sn-Pb) perovskite solar cells (PSCs) preserves
84.7\% efficiency after 450 h of maximum power point tracking (MPPT)
(champion 24.1\%) {[}41{]}, and doping with 0.8 mol\% fluorinated
methylammonium iodide (FMAI) maintains operational stability over 2000 h
(champion 25.67\%) {[}42{]}. Extended defect formation presents an
associated thermodynamic dilemma: dimethylammonium (DMA+) phase
engineering suppresses Ruddlesden-Popper antiphase boundary (RP-APB)
defects in beta-(DMA,Cs)PbI3 without nucleating the photo-inactive delta
phase from edge-sharing {[}PbI6{]}4- octahedra (champion 20.23\%)
{[}43{]}. Defect healing also follows an explicit kinetic hierarchy:
ordered anion-cation modification (OACM) demonstrates that anionic
vacancies must be repaired before cation passivation to match bromide
diffusion kinetics, enabling a 1.83 eV wide-bandgap (WBG) cell with an
open-circuit voltage (Voc) of 1.33 V and 19.21\% stabilized efficiency
(champion 26.2\%) {[}44{]}.

At device contacts, dimensional heterostructures and buried molecular
interlayers block defect propagation and mobile ion accumulation
{[}45,46,47,48,49,50,51{]}. A tilted one-dimensional (1D) on
three-dimensional (3D) heterointerface directed by a polyfluorinated
ammonium iodide regulator preserves vertical transport while physically
obstructing defect propagation and ion migration (champion 24.94\% PCE)
{[}45{]}. In two-dimensional (2D) capping layers, octylammonium bromide
(OABr) concentration dictates phase shifts toward higher n-values
(champion 19.29\% PCE) {[}48{]}, whereas chemical polishing removes
unreacted formamidinium iodide (FAI) and residual lead iodide (PbI2) to
enforce vertical 2D orientation (champion 21.36\% PCE) {[}47{]}.
Solid-state halide exchange using bromide-rich nanocrystals passivates
iodide vacancies and suppresses ion drift (uncertified preprint champion
22.19\%) {[}46{]}. At buried electron-transport layer (ETL) interfaces,
2-formylbenzenesulfonic acid sodium salt (2-FAS) bridges self-assembled
monolayers (SAM) to perovskite (champion 26.21\% PCE, 86.15\% FF)
{[}49{]}, acesulfame potassium (Ace-K) chelates Sn4+ and oxygen
vacancies on tin oxide (SnO2) to eliminate voids, retaining 92.4\%
efficiency after 1000 h of MPPT on an uncertified 0.0729 cm2 cell
(champion 26.17\%) {[}50{]}, and transient lead thiocyanate (Pb(SCN)2)
vapor eliminates residual PbI2 traps (champion 25.2\%) {[}51{]}.

\subsection{5. Wide-Bandgap Absorbers and Tandem
Integration}\label{wide-bandgap-absorbers-and-tandem-integration}

Solution deposition of hole-selective self-assembled monolayers (SAMs)
on textured silicon bottom cells suffers from inhomogeneous
conformality, where uneven coverage of
{[}2-(3,6-dimethoxy-9H-carbazol-9-yl)ethyl{]}phosphonic acid (MeO-2PACz)
precipitates residual lead iodide (PbI2) and causes interfacial shunting
{[}7,52{]}. Cooperative self-assembled monolayer (SAM) strategies
incorporating co-adsorbents with a methoxy terminal group and flexible
alkyl chain improve substrate wettability and molecular packing on
pyramids, lifting the uncertified power conversion efficiency (PCE) of
two-terminal (2T) perovskite/silicon tandems with 1.68 eV wide-bandgap
(WBG) absorbers to 32.16\% {[}53{]}. To prevent degradation during
atomic layer deposition (ALD), sacrificial coordination suppresses SAM
aggregation via hydrogen bonding and shields the anchored monolayer,
enhancing adhesion (ASTM D3359) to achieve an operational T90 of 1265
hours and a certified PCE of 26.56\% (champion 27.03\%) on a miniature
0.062 cm2 aperture {[}54{]}. Inorganic hole-transport layer (HTL) stacks
offer an alternative to vulnerable organophosphonates: pairing nickel
oxide (NiOx) with dye N719 achieves 96.4\% PCE retention after one year
of outdoor maximum power point tracking (MPPT) {[}55{]}. In polymeric
HTLs, dilauroyl peroxide doping of
poly{[}bis(4-phenyl)(2,4,6-trimethylphenyl)amine{]} (PTAA) raises the
water contact angle from 60.18° to 80.39°, boosting champion efficiency
to 24.82\% in n-i-p and 26.04\% in p-i-n perovskite solar cells (PSCs)
{[}56{]}.

Concurrently, disordered nucleation, strain, and ion drift across
textured pyramids compromise WBG absorber stability {[}57,58{]}.
Pre-depositing a cesium chloride (CsCl) seed layer across uneven SAMs
homogenizes organic precursor incorporation and eliminates interfacial
PbI2, securing a certified PCE of 29.7\% (30.3\% champion) {[}7{]}.
Solvent-free sequential vapor evaporation avoids liquid dewetting
entirely; in situ X-ray diffraction (XRD) confirms accelerated lead
halide conversion on textured silicon over planar substrates, yielding a
stabilised 27.0\% PCE {[}59{]}. Post-deposition healing with
methylammonium thiocyanate (MASCN) induces secondary crystallization
into columnar grains spanning the full film thickness, enabling an
encapsulated 1.164 cm2 device to achieve a certified 30.77\% PCE and
retain initial performance over 3400 hours of continuous one-sun MPPT in
ambient conditions {[}3{]}. To counteract mechanical distortion,
two-dimensional (2D) perovskite seeding directs regular compressive
strain into the three-dimensional (3D) WBG lattice, delivering 20.83\%
PCE with an 84.72\% fill factor (FF) in single junctions and 29.47\% in
tandems {[}60{]}. In the bulk absorber, ion modulation elevates the ion
migration activation barrier, lifting uncertified champion tandem PCE to
33.02\% from 31.12\% for control devices {[}61{]}.

At the tandem interconnect, optical and resistive losses restrict device
operation. In perovskite-organic tandems, a zinc sulfide (ZnS)
interlayer converts isolated gold clusters into laterally planarized
domains, establishing low-loss recombination pathways at only
\textasciitilde77.4\% metallic coverage while suppressing localized
surface plasmon resonance to yield a 27.1\% champion PCE and a T80 of
1400 hours {[}6{]}. Translating 2T tandems beyond 33\%-34\% laboratory
benchmarks toward commercial modules requires overcoming texturing
shunts, near-infrared (NIR) parasitic absorption, balance-of-system
(BOS) costs, and the defect cascades that penalize large-area scaling
{[}57,58{]}.

\begin{figure}[htbp]
\centering
\includegraphics[width=0.95\textwidth]{fig/F3_effort_map.pdf}
\caption{Where the month's mechanistic effort was directed. Depth-tier papers by mechanism axis, stacked by device family.}
\end{figure}

\subsection{6. Operational Stability and Degradation
Pathways}\label{operational-stability-and-degradation-pathways}

Operational degradation in inverted perovskite solar cells (PSCs)
originates primarily from coupled interfacial redox reactions, lattice
strain transitions, and electrode interdiffusion. At the buried
hole-transport layer (HTL), acidic carbazole-based phosphonic acid
self-assembled monolayers (SAMs) activate residual dimethyl sulfoxide
(DMSO), driving iodide redox decomposition during wet processing
{[}4{]}. Under subsequent light operation, this interfacial degradation
cascades into the electron-transport layer (ETL): formamidinium iodide
(FAI) photo-oxidation generates iodine radicals that induce
phenyl-C61-butyric acid methyl ester (PCBM) dimerization, while
fullerene species accelerate FAI deprotonation in a self-reinforcing
photochemical loop {[}62{]}. Interfacial radical scavengers {[}62{]},
cross-linkable self-assembled monolayer (SAM) linker architectures
{[}63{]}, and dipole-balancing molecular co-assemblies {[}64,65,66{]}
mitigate these decomposition cascades. Nevertheless, stark discrepancies
between uncertified champion power conversion efficiency (PCE) values
and certified metrics, including a champion 27.06\% PCE dropping to a
certified 20.11\% on small 0.09 cm2 apertures {[}66{]}, reveal
unresolved interfacial instability and measurement drift.

Within the perovskite absorber, vertical lattice heterogeneity
accelerates lead and iodide vacancy formation under operating stress
{[}67{]}. Thermal and operational biases exacerbate internal lattice
strain, promoting polymorphic degradation from the photoactive alpha
phase to the non-photoactive delta phase; relieving this strain through
interstitial incorporation of 2 mol\% sodium {[}68{]} or templating via
two-dimensional (2D) intermediate phases {[}69{]} stabilizes the active
crystal lattice. At the rear contact, halide migration triggers severe
metal interdiffusion, requiring diffusion barriers such as 5 nm chromium
interlayers to prevent contact failure {[}70,71{]}.

A critical systemic weakness across these investigations is the severe
lack of standardized aging protocols. Only 5 of 144 surveyed papers
reported a time to 80\% initial performance (T80) value: 1400 h {[}6{]},
900 h {[}19{]}, 864 h {[}71{]}, 800 h {[}68{]}, and 720 h {[}72{]}.
Concurrently, only 2 studies adhered to formal International Summit on
Organic Photovoltaic Stability (ISOS) protocols, specifically ISOS-L-1
{[}64{]} and ISOS-L-2 {[}62{]}. Most stability tests rely on
unencapsulated continuous illumination without standardized maximum
power point tracking (MPPT) {[}73{]}. Consequently, existing literature
primarily probes short-term photochemical burn-in, obscuring critical
field degradation pathways including reverse-bias breakdown,
moisture-accelerated ion migration, and thermo-mechanical fatigue under
diurnal operational cycling.

\begin{figure}[htbp]
\centering
\includegraphics[width=0.95\textwidth]{fig/F4_stability_evidence.pdf}
\caption{The month's operational-stability evidence in full. Left, every reported time to 80\% of initial performance. Right, the stability protocol labels used.}
\end{figure}

\subsection{7. Scale-Up and the Area
Penalty}\label{scale-up-and-the-area-penalty}

The scaling of perovskite solar cells (PSCs) from laboratory coupons to
large-area modules incurs an efficiency penalty dominated by spatial
mass-transfer disparities, stochastic nucleation, and resistive or
environmental losses. Non-uniform drying across decimeter areas induces
pinholes and phase segregation; replacing vertical diffusion-limited
solvent evaporation with convective or engineered antisolvent extraction
suppresses these defects {[}5,11,74{]}. Closed-loop slit-flushing
antisolvent deposition imposes horizontal forced convection over 100 ×
100 mm substrates to achieve a certified power conversion efficiency
(PCE) of 17.35\% on a 61.2 cm2 module {[}5{]}. Static antisolvent
bathing with isopropanol admixtures homogenizes crystallization into
larger crystalline nanodomains verified by grazing-incidence wide-angle
X-ray scattering (GIWAXS), enabling an uncertified 27.7\% champion PCE
on a 25 cm2 two-terminal (2T) tandem device {[}11{]}, while polar
additives in slot-die precursors yield 30 nm films with 1.86 nm
roughness over 100 cm2 {[}74{]}. In parallel, chemical coordination and
heterogeneous seeding regulate nucleation barriers to prevent premature
precipitation {[}12,75,76{]}. Polyfluoroarene additives elevate
effective nucleation barriers to convert rapid aggregation into uniform
film growth, achieving an uncertified module champion PCE of 27.01\% on
58.51 cm2 with small-area non-radiative voltage losses of 55 mV
{[}12{]}. Purine tautomers alter lead-nitrogen coordination and hydrogen
bonding, confirmed by density functional theory (DFT), yielding a
certified 26.51\% PCE (champion 27.05\%) at 12.95 cm2 {[}75{]}, whereas
cesium lead triiodide (CsPbI3) nanocube seeding enforces controlled
crystallization and depth-dependent Cs+ incorporation in slot-die coated
formamidinium lead triiodide (FAPbI3) {[}76{]}. Beyond nucleation,
solid-state transformation barriers and resistive contacts limit module
performance {[}77,78{]}. Laser-structured channels in lead iodide (PbI2)
provide deterministic diffusion pathways that lower phase transformation
activation barriers, reaching an uncertified 22.83\% PCE across 22.95
cm2 modules {[}77{]}. For charge extraction, cerium iodide (CeI3)
interfacial modifiers establish inorganic hole contacts, yielding
25.19\% PCE in 0.113 cm2 cells and 23.37\% in 21 cm2 modules with a
98.3\% geometric fill factor (GFF) {[}78{]}. Finally, scalable slot-die
processing remains constrained by environmental sensitivity and
mechanical degradation {[}79,80{]}. Ambient humidity induces marked PCE
degradation when exceeding 30\% relative humidity (RH) {[}79{]}, yet
incorporating Ti3C2Tx MXene nanospheres extends continuous illumination
stability from 400 h to over 1700 h in all-slot-die-coated PSCs, albeit
at an uncertified 18.2\% champion PCE {[}80{]}.

\begin{figure}[htbp]
\centering
\includegraphics[width=0.95\textwidth]{fig/F2_area_penalty.pdf}
\caption{Efficiency against reported active or aperture area. The shaded band marks the sub-0.3 cm$^2$ laboratory-cell regime in which most champion values are measured.}
\end{figure}

\subsection{8. Research Gaps and
Outlook}\label{research-gaps-and-outlook}

The performance divergence between single junctions and tandems defines
the first gap: the best certified perovskite/silicon tandem this month
was 33.66\% power conversion efficiency (PCE) {[}1{]}, while certified
single junctions clustered near 27\%, showing that the wide-bandgap
(WBG) subcell limiting the tandem is not the one most papers optimise.
While chemical strategies claim to homogenize mixed halides in pristine
films {[}81,82{]}, whether halide homogenisation in wide-bandgap
absorbers suppresses phase segregation under operating bias remains
unproven because photocarrier gradients trigger iodide clustering.
Operando spatially resolved photoluminescence (PL) spectroscopy under
continuous maximum power point tracking (MPPT) load at elevated
temperatures would resolve whether compositional uniformity survives
operational electric fields.

A second gap concerns buried-interface chemistry that survives thermal
cycling rather than only lifting initial efficiency. Molecules and
self-assembled monolayers (SAMs) passivate buried defects to lift
open-circuit voltage (Voc) on pristine substrates {[}63,66{]}, but
mismatched thermal expansion coefficients generate interfacial shear
strain that induces delamination. Only 5 of 144 closely read papers
reported a time to 80\% of initial performance (T80) lifetime, the
longest 1400 h {[}6{]}, short against a 25-year module warranty. In situ
temperature-dependent adhesion testing with impedance spectroscopy
during thermal cycling would reveal whether interfacial bonding survives
thermomechanical stress.

A third gap lies in ion-migration timescales measured under load rather
than in the dark. Passivation schemes report reduced ion drift
{[}44,61{]}, yet extracting activation energies in the dark ignores how
photogenerated carriers screen internal fields and accelerate vacancy
migration under illumination. Operando transient ion-drift measurements
under illumination during continuous MPPT would resolve ionic transit
timescales and defect accumulation under load.

The fourth gap is whether passivation that lifts Voc on a 0.1 cm2 cell
survives slot-die coating at module area. Only 27 of 531 works addressed
scale-up at all, and the largest reported aperture was 61.2 cm2 {[}5{]},
orders of magnitude below a commercial module. Additives that distribute
evenly during spin coating aggregate under meniscus shear, creating
non-uniform trap landscapes {[}12,76{]}. Operando PL imaging during
slot-die coating would track drying kinetics and isolate spatial Voc
losses. Standardizing stability assessments to International Summit on
Organic Photovoltaic Stability (ISOS) protocols under operational load
and validating passivation across large-area deposition would make next
month's literature considerably more useful.

\subsection{Acknowledgements}\label{acknowledgements}

The authors would like to acknowledge the financial support from Xiamen
University Malaysia through the Xiamen University Malaysia Research Fund
(XMUMRF/2022-C10/IENG/0046), and from National Natural Science
Foundation of China (Ref. No.: 52503390).

\subsection{Declaration of Competing
Interest}\label{declaration-of-competing-interest}

The authors declare no conflict of interest.

\subsection{Data Availability
Statement}\label{data-availability-statement}

The corpus table for the month, the per-paper extraction records with
their verbatim source quotations, the statistics file underlying every
number in the text, and the gate report are provided as Supplementary
Information. Abstract text and full texts are not redistributed.

\subsection{AI Usage Declaration}\label{ai-usage-declaration}

This manuscript was prepared as part of our Research Project for making
a reliable Autonomous Researcher Agent. The agent executed the monthly
review pipeline without human intervention: it harvested the indexed
literature for the month through the OpenAlex, Semantic Scholar,
Crossref and arXiv interfaces, applied a deterministic topical filter,
ranked and selected the papers read closely, extracted every
quantitative claim into a structured record binding each number to a
verbatim quotation from that paper's own abstract, computed the
statistics and figures from those records, and drafted this text from
the resulting evidence briefs. Literature harvesting, selection,
statistics, figures and gate checking are deterministic scripts rather
than model output. Structured extraction used qwen3.8-flash through the
opencode interface; the prose was written by gemini-3.8-flash-high
through the agy interface; independent review used opus at high
reasoning effort. No figure is a model-generated image: every panel is
plotted directly from the extracted records. Every quantitative
statement in the main text resolves to the machine-generated statistics
file shipped as Supplementary Information, and every citation resolves
to a record carrying its source quotation. The author reviewed and
edited the content, verified the selection, and takes full
responsibility for the content of the published article.

\subsection{References}\label{references}

\begin{enumerate}
\def\labelenumi{\arabic{enumi}.}
\tightlist
\item
  Zeng, L. et al.~Non-competitive additives regulate molecular assembly
  for compact monolayers in perovskite/silicon tandems. \emph{Nature
  Communications} (2026).
  \url{https://doi.org/10.1038/s41467-026-76266-0}
\item
  Zhang, D. et al.~Thienyl Self‐Assembled Molecule Enabled Stable
  Perovskite/TOPCon Silicon Tandem. \emph{Angewandte Chemie
  International Edition} (2026).
  \url{https://doi.org/10.1002/anie.5241611}
\item
  Kong, W. et al.~Healing intervention for improving the efficiency and
  stability of tandem devices on industrial textured silicon.
  \emph{Nature Communications} (2026).
  \url{https://doi.org/10.1038/s41467-026-76713-y}
\item
  Liang, Z. et al.~Redirecting wet--interfacial redox pathways for
  efficient inverted perovskite solar cells. \emph{Science} (2026).
  \url{https://doi.org/10.1126/science.aeg8415}
\item
  Xu, Y. et al.~Closed‐Loop Slit‐Flushing Antisolvent Strategy Enables
  Controlled Three‐Dimensional Mass Transfer for Large‐Area Perovskite
  Solar Cells. \emph{Small} (2026).
  \url{https://doi.org/10.1002/smll.75068}
\item
  Ren, H. et al.~Pushing the Interconnection Limit in Perovskite‐Organic
  Tandem Solar Cells Through Recombination Channels Delocalization.
  \emph{Angewandte Chemie International Edition} (2026).
  \url{https://doi.org/10.1002/anie.6548693}
\item
  Škorjanc, V. et al.~CsCl seed layer homogenizes co-evaporated
  perovskite growth for high-efficiency, fully textured
  perovskite-silicon tandem solar cells. \emph{Joule} (2026).
  \url{https://doi.org/10.1016/j.joule.2026.102645}
\item
  Xie, Z. et al.~Stacking-order-driven interfacial dipole flip and
  chlorine-iodine exchange for high-performance perovskite-organic
  tandem solar cells. \emph{Nature Communications} (2026).
  \url{https://doi.org/10.1038/s41467-026-76786-9}
\item
  Pan, Y. et al.~Self-Assembled Multilayers Reduce Interfacial Energy
  Loss in Perovskite Solar Cells. \emph{Nature Communications} (2026).
  \url{https://doi.org/10.1038/s41467-026-76497-1}
\item
  Feng, Y. et al.~Ga Doping Enables (110) Facet Control and Surface
  Passivation of SnO 2 for Efficient and Stable Inverted Perovskite
  Solar Cells. \emph{Angewandte Chemie International Edition} (2026).
  \url{https://doi.org/10.1002/anie.3095986}
\item
  Pallotta, R. et al.~Static antisolvent processing for scalable and
  sustainable perovskite photovoltaics. \emph{Communications Materials}
  (2026). \url{https://doi.org/10.1038/s43246-026-01321-8}
\item
  Zhang, T. et al.~Synergistic Crystallization Regulation by Substituent
  Electronic Effects of Polyfluoroarenes Toward Large Area Perovskite
  Solar Module Fabrication. \emph{Advanced Materials} (2026).
  \url{https://doi.org/10.1002/adma.74469}
\item
  Qi, J. et al.~Geometrically Asymmetric Phosphonium Zwitterions Enable
  Efficient Printable Mesoscopic Perovskite Solar Cells. \emph{Advanced
  Materials} (2026). \url{https://doi.org/10.1002/adma.74629}
\item
  Vaitukaitytė, D. et al.~Anchor Engineering of Hole‐Selective
  Enamine‐Based Self‐Assembling Monolayers for Perovskite Solar Cells.
  \emph{Advanced Science} (2026).
  \url{https://doi.org/10.1002/advs.76950}
\item
  Feng, X. et al.~Designing Efficient Inverted Perovskite Solar Cells
  with Self-Assembled Monolayer Hole Transport Layers. \emph{Nano-Micro
  Letters} (2026). \url{https://doi.org/10.1007/s40820-026-02327-0}
\item
  Liu, H. et al.~Hidden but Not Forgotten: Unraveling and Engineering
  the Buried Interface of Inverted Perovskite Solar Cells. \emph{Solar
  RRL} (2026). \url{https://doi.org/10.1002/solr.70473}
\item
  Ge, Y. et al.~Interfacial Engineering Strategies of Self-Assembled
  Monolayers for Inverted Perovskite Solar Cells. \emph{Nanomaterials}
  (2026). \url{https://doi.org/10.3390/nano16160989}
\item
  Ding, S. et al.~Inhibiting Buried Mechanical Failure via Vapor‐Induced
  Chemical Reconstruction for Durable Solar Cells. \emph{Small} (2026).
  \url{https://doi.org/10.1002/smll.75122}
\item
  Wang, X. et al.~Clean Tin Oxide via Volatile Additive Engineering for
  High‐performance Perovskite Photovoltaics. \emph{Advanced Energy
  Materials} (2026). \url{https://doi.org/10.1002/aenm.71514}
\item
  Sun, Y. et al.~Synergistic Passivation Effect via a Bidentate
  Anchoring Molecule between the Perovskite Layer and Electron Transport
  Layer for Perovskite Solar Cells. \emph{ACS Applied Energy Materials}
  (2026). \url{https://doi.org/10.1021/acsaem.6c01436}
\item
  Xie, F. et al.~Buried Interface Modulation via Molecular Dipole
  Passivation for High‐Efficiency Methylammonium‐Free Pb--Sn Perovskite
  Solar Cells. \emph{Small} (2026).
  \url{https://doi.org/10.1002/smll.75089}
\item
  Zhu, Y. et al.~Bifunctional NiO x /Perovskite Interface Engineering
  for Efficient Air-Processed Wide-Bandgap Perovskite Solar Cells.
  \emph{ACS Applied Energy Materials} (2026).
  \url{https://doi.org/10.1021/acsaem.6c01708}
\item
  Li, Y. et al.~Structural Compensation Strategy: Composite NiO x Hole
  Transport Layers Enabling High‐Efficiency and Scalable Inverted
  Perovskite Solar Cells. \emph{Small} (2026).
  \url{https://doi.org/10.1002/smll.75056}
\item
  Hu, C. et al.~Hydrogen-bond network-reinforced self-assembled
  monolayer enables high efficiency and operational stability in
  pure-iodide wide-bandgap perovskite solar cells. \emph{DeCarbon}
  (2026). \url{https://doi.org/10.1016/j.decarb.2026.100190}
\item
  Wang, H. et al.~Dual Modulation of the Buried Interface in Inverted
  Perovskite Solar Cells With Heteroatom‐Conjugated Asymmetric
  Self‐Assembled Molecules. \emph{Advanced Functional Materials} (2026).
  \url{https://doi.org/10.1002/adfm.77712}
\item
  Qin, X. et al.~Molecular Regulation toward SAM/Perovskite Interface
  for Efficient p‐i‐n Perovskite Solar Cells. \emph{Advanced Energy
  Materials} (2026). \url{https://doi.org/10.1002/aenm.71482}
\item
  Zhang, W. et al.~Dipole-engineered co-adsorbed self-assembled
  monolayers boost Cs-based inorganic Pb-Sn perovskite tandem efficiency
  beyond 20\%. \emph{Nature Communications} (2026).
  \url{https://doi.org/10.1038/s41467-026-76847-z}
\item
  Zeng, X. et al.~Dipole-Moment-Tuned Self-Assembled Monolayers for
  Efficient Perovskite Solar Cells. \emph{ACS Applied Materials \&
  Interfaces} (2026). \url{https://doi.org/10.1021/acsami.6c11467}
\item
  Zhou, S. et al.~The Impact of C60-Self-Assembled Monolayer
  Electron-Transport Layers in Negative--Intrinsic--Positive Perovskite
  Solar Cells. \emph{ACS Energy Letters} (2026).
  \url{https://doi.org/10.1021/acsenergylett.6c01485}
\item
  Yuan, S. et al.~Molecularly Templated Buried Interfaces for Inverted
  Perovskite Solar Cells. \emph{Angewandte Chemie International Edition}
  (2026). \url{https://doi.org/10.1002/anie.6932848}
\item
  Lu, C. et al.~Near-Surface Sulfidation-Induced p-n Junction Enabling
  High-Efficiency Inverted Perovskite Solar Cells. \emph{ACS Applied
  Materials \& Interfaces} (2026).
  \url{https://doi.org/10.1021/acsami.6c07884}
\item
  Yang, J. et al.~Efficient and Stable Inverted Perovskite Solar Cells
  Via a Multi‐Arm Donor--Acceptor Dipole Molecular Bridge.
  \emph{Angewandte Chemie International Edition} (2026).
  \url{https://doi.org/10.1002/anie.3716107}
\item
  Kang, R. et al.~Thiourea Derivatives for Mitigating
  Photodegradation-Mediated Lead Leakage in Sustainable Perovskite Solar
  Cells and Modules. \emph{ACS Energy Letters} (2026).
  \url{https://doi.org/10.1021/acsenergylett.6c02033}
\item
  Huang, J. et al.~Sustainable Stabilization of Perovskite Solar Cells
  Using Multifunctional Ammonium Alginate for over 26\% Efficiency.
  \emph{Advanced Functional Materials} (2026).
  \url{https://doi.org/10.1002/adfm.78067}
\item
  Wang, C.-. et al.~Sulfamethoxazole-Induced Synergistic Interfacial
  Passivation for Efficient and Stable Perovskite Solar Cells. \emph{The
  Journal of Physical Chemistry Letters} (2026).
  \url{https://doi.org/10.1021/acs.jpclett.6c02253}
\item
  Chen, L. et al.~Dual‐Action Sulfonium‐Amino Acid Additive Enabling
  Crystallization Regulation and Defect Passivation for High‐Performance
  Carbon‐Electrode Printable Mesoscopic Perovskite Solar Cells.
  \emph{Angewandte Chemie International Edition} (2026).
  \url{https://doi.org/10.1002/anie.5667084}
\item
  Liu, X. et al.~Novel Multi-Pyridyl-Functionalized Spirobifluorenes as
  Precursor Additives for Improving Efficient and Stable Perovskite
  Solar Cells. \emph{ACS Applied Materials \& Interfaces} (2026).
  \url{https://doi.org/10.1021/acsami.6c13903}
\item
  Du, C. et al.~Synergistic Dipole and Dual‐Passivation Enabled Highly
  Efficient Two‐Dimensional Perovskite Solar Cells. \emph{Advanced
  Functional Materials} (2026). \url{https://doi.org/10.1002/adfm.77504}
\item
  Liang, J. et al.~Solid-solution-induced lattice hardening reinforces
  perovskite stiffness for stable solar cell. \emph{National Science
  Review} (2026). \url{https://doi.org/10.1093/nsr/nwag484}
\item
  Ali, W. et al.~Suppressing Non‐Radiative Recombination in Perovskite
  Solar Cells Through Integrated Tuning of Out‐of‐Plane Lattice Texture
  and Strain. \emph{Advanced Energy Materials} (2026).
  \url{https://doi.org/10.1002/aenm.71474}
\item
  Zhao, J. et al.~Dynamic Strain Homogenization Regulation for
  Thermomechanically Robust Sn─Pb Perovskite Solar Cells. \emph{Advanced
  Energy Materials} (2026). \url{https://doi.org/10.1002/aenm.71497}
\item
  Mikuciunaite, J. et al.~Fluorinated Methylammonium Cation Containing
  Perovskite Solar Cells With Over 25\% Power Conversion Efficiency.
  \emph{Advanced Energy Materials} (2026).
  \url{https://doi.org/10.1002/aenm.71396}
\item
  Yin, Z. et al.~Antiphase Boundaries Regulate Phase Stability and
  Performance in DMA + ‐Assisted CsPbI 3 ‐Based Perovskites.
  \emph{Advanced Materials} (2026).
  \url{https://doi.org/10.1002/adma.74538}
\item
  Ouyang, D. et al.~Diminishing Defect and Ion Migration in
  Perovskite‐Organic Tandem Solar Cells by Ordered Anion‐Cation
  Modification Approach. \emph{Advanced Functional Materials} (2026).
  \url{https://doi.org/10.1002/adfm.77855}
\item
  Chu, H. et al.~Polyfluorinated‐Regulator‐Assisted Scaffold‐Directed
  Tilted 1D/3D Heterointerfaces for High‐Voltage and Stable Perovskite
  Solar Cells. \emph{Angewandte Chemie International Edition} (2026).
  \url{https://doi.org/10.1002/anie.8256065}
\item
  Rout, P. et al.~Halide Exchange via Perovskite Nanocrystals for
  Interfacial Defect Passivation in Solar Cells. \emph{ChemRxiv} (2026).
  \url{https://doi.org/10.26434/chemrxiv.15007814/v1}
\item
  Yang, F. et al.~Constructing Vertically Oriented 2D Perovskite
  Passivation Layer by Surface Chemical Polishing for Efficient
  Carbon‐Based Perovskite Solar Cells. \emph{Small} (2026).
  \url{https://doi.org/10.1002/smll.75536}
\item
  Yang, Y. et al.~Precise Management of n‐Value Distribution at the
  2D/3D Interface Enabling Efficient Wide‐Bandgap Perovskite Solar
  Cells. \emph{Advanced Science} (2026).
  \url{https://doi.org/10.1002/advs.76791}
\item
  Zha, L. et al.~A Buried Sulfonate Molecular Bridge for Synchronous
  Charge Transport and Defect Passivation in High-Performance Perovskite
  Solar Cells. \emph{Nano-Micro Letters} (2026).
  \url{https://doi.org/10.1007/s40820-026-02278-6}
\item
  Sun, W. et al.~Competitive Ligand Modulation Reduces Buried
  Interfacial Voids and Defects Enabling 26.17\% Efficient Perovskite
  Solar Cells. \emph{Small} (2026).
  \url{https://doi.org/10.1002/smll.75360}
\item
  Chen, W. et al.~Lowering the δ-to-α Phase-Transition Barrier via a
  Transient Vapor Mediator in Perovskite Formation. \emph{Journal of the
  American Chemical Society} (2026).
  \url{https://doi.org/10.1021/jacs.6c12975}
\item
  Park, S.J. et al.~Self-Assembled Monolayers for High-Efficiency
  Perovskite/Si Tandem Solar Cells. \emph{ACS Energy Letters} (2026).
  \url{https://doi.org/10.1021/acsenergylett.6c01966}
\item
  Gong, H. et al.~Rationally Designed Cooperative Self‐Assembled
  Monolayers Enable Holistic Modulations of Wide‐Bandgap Perovskites for
  High‐Performance Perovskite‐Silicon Tandem Solar Cells. \emph{Advanced
  Energy Materials} (2026). \url{https://doi.org/10.1002/aenm.71512}
\item
  Yuan, J. et al.~Universal Sacrificial Coordination Strategy for
  ALD‐Resilient SAMs Achieving High‐Performance Perovskite/Organic
  Tandem Solar Cells. \emph{Advanced Materials} (2026).
  \url{https://doi.org/10.1002/adma.74461}
\item
  Bristow, H. et al.~Material Selection through Accelerated Aging
  Enables \textgreater95\% Performance Retention of a
  Perovskite--Silicon Tandem Solar Cell after 1 Year of Outdoor
  Monitoring. \emph{ACS Energy Letters} (2026).
  \url{https://doi.org/10.1021/acsenergylett.6c01363}
\item
  Lan, Z. et al.~Hydrophobic Dilauroyl Peroxide Improves the Electrical
  Performance and Stability of PTAA in both n-i-p and p-i-n Perovskite
  Solar Cells. \emph{The Journal of Physical Chemistry Letters} (2026).
  \url{https://doi.org/10.1021/acs.jpclett.6c02136}
\item
  Li, L. et al.~Large‐Scale Perovskite Photovoltaics: Mitigating Defect
  Challenges in Upscaled Modules and Monolithic Tandem Cells.
  \emph{Advanced Functional Materials} (2026).
  \url{https://doi.org/10.1002/adfm.202532038}
\item
  Uddin, K.B. Recent Advances, Efficiency Limits, and Commercialization
  Challenges in Perovskite-Based Tandem Solar Cells: A Comprehensive
  Review. \emph{Zenodo (CERN European Organization for Nuclear
  Research)} (2026). \url{https://doi.org/10.5281/zenodo.21790600}
\item
  Mahmoud, M.A.A. et al.~Impact of silicon substrate topography on
  sequentially evaporated perovskite film growth. \emph{Joule} (2026).
  \url{https://doi.org/10.1016/j.joule.2026.102646}
\item
  Zhang, T. et al.~Dimensional Seed Engineering of Strain in
  Wide-Bandgap Perovskites for Efficient and Stable Tandem Solar Cells.
  \emph{ACS Nano} (2026). \url{https://doi.org/10.1021/acsnano.6c07355}
\item
  Zhao, Z. et al.~High‐Performance Perovskite/Silicon Tandem Solar Cells
  via Ion Modulation. \emph{Advanced Functional Materials} (2026).
  \url{https://doi.org/10.1002/adfm.77637}
\item
  Zhang, J. et al.~Decoupling Bidirectional Photochemical Degradation
  via Radical Scavenging for Stable Inverted Perovskite Solar Cells.
  \emph{Angewandte Chemie International Edition} (2026).
  \url{https://doi.org/10.1002/anie.4221605}
\item
  Yan, X. et al.~Self-assembled monolayers with cross-linkable header on
  the linker for durable perovskite photovoltaics. \emph{Science
  Advances} (2026). \url{https://doi.org/10.1126/sciadv.aec3096}
\item
  Luo, C. et al.~Symmetry-breaking co-assembly of conjugated molecules
  boosts perovskite photovoltaics. \emph{Nature Communications} (2026).
  \url{https://doi.org/10.1038/s41467-026-76848-y}
\item
  Guo, Y. et al.~Regulating Molecular Adsorption Configuration via
  Spatially Balanced Polarity for High‐Performance Inverted Perovskite
  Solar Cells. \emph{Advanced Functional Materials} (2026).
  \url{https://doi.org/10.1002/adfm.77942}
\item
  Wang, K. et al.~Acid‐Base Complexation Induced Dipole Engineering for
  Durable Inverted Perovskite Photovoltaics. \emph{Advanced Materials}
  (2026). \url{https://doi.org/10.1002/adma.74438}
\item
  Abbas, Z. et al.~Suppression of Lattice Heterogeneity for Efficient
  Inverted Perovskite Solar Cells. \emph{Advanced Functional Materials}
  (2026). \url{https://doi.org/10.1002/adfm.77729}
\item
  Tang, Z. et al.~Mitigating Stress‐Induced Nonphotoactive Phase
  Transition Through Sodium Sulfonate Engineering for Stable and
  Efficient Perovskite Solar Cells. \emph{Angewandte Chemie
  International Edition} (2026).
  \url{https://doi.org/10.1002/anie.2425823}
\item
  Ma, X. et al.~Regulating lattice evolution via an ordered
  two-dimensional intermediate phase for resilient perovskite solar
  cells. \emph{Nature Communications} (2026).
  \url{https://doi.org/10.1038/s41467-026-76425-3}
\item
  Lukas, T. et al.~Beyond the Gold Standard: Towards Industrially Viable
  Electrodes for Durable Perovskite Solar Cells. \emph{Advanced Energy
  Materials} (2026). \url{https://doi.org/10.1002/aenm.71372}
\item
  Liu, N. et al.~Two Birds With One Stone: Modifier Interfacial
  Engineering for High‐Efficiency Stable p‐i‐n All‐Inorganic
  Perovskite/Organic Tandem Solar Cells. \emph{Advanced Functional
  Materials} (2026). \url{https://doi.org/10.1002/adfm.77805}
\item
  Lin, S. et al.~Steric Hindrance‐Assisted Coordination for Efficient
  and Stable Tin‐Based Perovskite Solar Cells. \emph{Solar RRL} (2026).
  \url{https://doi.org/10.1002/solr.70454}
\item
  Wang, S. et al.~Graded Supramolecular Frameworks for Extending
  Space‐Charge Modulation Regions and Boosting Carrier Dynamics in
  Perovskite Solar Cells. \emph{Advanced Materials} (2026).
  \url{https://doi.org/10.1002/adma.74489}
\item
  Jin, D. et al.~Slot-Die Coating 30 nm Thick Perovskite Films for
  Large-Area Red Light-Emitting Diodes. \emph{The Journal of Physical
  Chemistry Letters} (2026).
  \url{https://doi.org/10.1021/acs.jpclett.6c02178}
\item
  Shi, J. et al.~Harnessing Molecular Tautomerism to Tune Precursor
  Coordination Chemistry for Efficient and Stable Perovskite Solar
  Cells. \emph{Angewandte Chemie International Edition} (2026).
  \url{https://doi.org/10.1002/anie.7848549}
\item
  Buyan‐Arivjikh, A. et al.~Overcoming Uncontrolled Nucleation and
  Cesium Salt Solubility Limits in Scalable Perovskite Deposition
  through All-Inorganic Perovskite Nanocrystal Seeding. \emph{ACS
  Applied Materials \& Interfaces} (2026).
  \url{https://doi.org/10.1021/acsami.6c11722}
\item
  Du, K. et al.~Laser‐Defined Reaction Topology Enables Controllable
  Solid‐State Transformations for Scalable Perovskite Photovoltaics.
  \emph{Advanced Materials} (2026).
  \url{https://doi.org/10.1002/adma.74828}
\item
  Zhang, H. et al.~Fully Inorganic Hole Contacts Enabled by CeI3 for
  High-Performance Perovskite Solar Modules. \emph{ACS Nano} (2026).
  \url{https://doi.org/10.1021/acsnano.5c21962}
\item
  Bonatto, A.E.T. et al.~Identifying an Ambient Humidity Processing
  Window for the Slot-Die Coating of Perovskite Solar Cells. \emph{ACS
  Applied Energy Materials} (2026).
  \url{https://doi.org/10.1021/acsaem.6c02083}
\item
  Le, S. et al.~Stabilization of All-Slot-Die-Coated Perovskite Solar
  Cells by MXene Nanospheres Unleashes Thousand-Hour Operational
  Lifetimes. \emph{ACS Applied Materials \& Interfaces} (2026).
  \url{https://doi.org/10.1021/acsami.6c07537}
\item
  Bi, L. et al.~Crystallization Modulation for Stable Wide‐Bandgap
  Perovskite Solar Cells and Modules. \emph{Advanced Materials} (2026).
  \url{https://doi.org/10.1002/adma.74388}
\item
  Gao, Z. et al.~Selenocyanate Driven Organic Cations Homogenization in
  Wide‐Bandgap Perovskite. \emph{Advanced Energy Materials} (2026).
  \url{https://doi.org/10.1002/aenm.71462}
\end{enumerate}

\end{document}
